% !TEX program = xelatex
\documentclass{英语试卷}
\begin{document}
\examtitle{山东省济宁、枣庄两地市 2025 年 4 月高考模拟试题}{英语}
\par 山东省济宁、枣庄两地市 2025 年 4 月高考模拟试题\par
\par 高三英语\par
\par 命题\&审题：山东省济宁市教科院\par
\begin{notice}
\par 1. 答卷前，考生务必将自己的姓名、准考证号填写在答题卡上。\par
\par 2. 回答选择题时，选出每小题答案后，用铅笔把答题卡上对应题目的答案标号涂黑。如需改动，用橡皮擦干净后，再选涂其他答案标号。回答非选择题时，将答案写在答题卡上，写在本试卷上无效。\par
\par 3. 考试结束后，将本试卷和答题卡一并交回。\par
\end{notice}
\parttitle{第一部分 听力（共两节，满分 30 分）}
\par 做题时，先将答案标在试卷上。录音内容结束后，你将有两分钟的时间将试卷上的答案转涂到答题卡上。\par
\sectitle{第一节（共5小题：每小题1.5分，满分7.5分）}
\par 听下面 5 段对话。每段对话后有一个小题，从题中所给的 A、B、C 三个选项中选出最佳选项。听完每段对话后，你都有 10 秒钟的时间来回答有关小题和阅读下一小题。每段对话仅读一遍。\par
\question{1. How does the woman sound?}
\begin{choices}(3)
  \item Confused.
  \item Disappointed.
  \item Anxious.
\end{choices}
\question{2. What does the woman ask about?}
\begin{choices}(3)
  \item A traveling route.
  \item A storage facility.
  \item A building entrance.
\end{choices}
\question{3. What is the man's problem?}
\begin{choices}(3)
  \item He is overcharged.
  \item He can't find the key.
  \item He is late for the checkout.
\end{choices}
\question{4. What do we know about the pumpkin milky tea?}
\begin{choices}(3)
  \item It is not fresh enough.
  \item It is unavailable now.
  \item It is a winter drink.
\end{choices}
\question{5. What are the speakers mainly talking about?}
\begin{choices}(3)
  \item A job application.
  \item A job change.
  \item A job interview.
\end{choices}
\sectitle{第二节（共15小题：每小题1.5分，满分22.5分）}
\par 听下面 5 段对话或独白。每段对话或独白后有几个小题，从题中所给的 A、B、C 三个选项中选出最佳选项。听每段对话或独白前，你将有时间阅读各个小题，每小题 5 秒钟；听完后，各小题将给出 5 秒钟的作答时间。每段对话或独白读两遍。 听第 6 段材料，回答第 6、7 题，\par
\question{6. What is the woman doing now?}
\begin{choices}(3)
  \item Repairing a road.
  \item Shopping in a store.
  \item Asking for directions.
\end{choices}
\question{7. How does the man suggest the woman return?}
\begin{choices}(3)
  \item Take bus Number 22.
  \item Take the underground.
  \item Take bus Number 7.
\end{choices}
\par 听第 7 段材料, 回答第 8 至 10 题,\par
\question{8. Where are the speakers probably?}
\begin{choices}(3)
  \item On a road.
  \item At a car service center.
  \item In a parking lot.
\end{choices}
\question{9. Why did the woman schedule her appointment for today?}
\begin{choices}(3)
  \item She completed some preparations yesterday.
  \item She doesn't have to work this afternoon.
  \item She will need the car tomorrow.
\end{choices}
\question{10. What does the man recommend doing at last?}
\begin{choices}(3)
  \item Postponing a service.
  \item Updating some records.
  \item Rescheduling an appointment.
\end{choices}
\par 听第 8 段材料,回答第 11 至 13 题。\par
\question{11. What causes the time conflict for the man?}
\begin{choices}(3)
  \item A medical appointment.
  \item A planned meeting.
  \item A college course.
\end{choices}
\question{12. Who will most probably take the man's shift?}
\begin{choices}(3)
  \item The woman.
  \item Ruby.
  \item Mr. Romero.
\end{choices}
\question{13. What does the woman advise the man to do when the manager arrives?}
\begin{choices}(3)
  \item Make a reminder note.
  \item Get some contact details.
  \item Report the progress of a task.
\end{choices}
\par 听第 9 段材料,回答第 14 至 16 题。\par
\question{14. What is the purpose of the charity marathon?}
\begin{choices}(3)
  \item To raise money for company's development.
  \item To provide training for volunteers.
  \item To support school education.
\end{choices}
\question{15. What voluntary work will the man choose to do?}
\begin{choices}(3)
  \item Managing the finish line.
  \item Helping with the set-up.
  \item Handing out water.
\end{choices}
\question{16. When does the race start?}
\begin{choices}(3)
  \item At 6 a. m.
  \item At 6:30 a. m.
  \item At 7 a. m.
\end{choices}
\par 听第 10 段材料, 回答第 17 至 20 题。\par
\question{17. What can visitors do along the Sarawak River?}
\begin{choices}(1)
  \item Visit handcraft stores.
  \item View the Astana palace.
  \item Explore side streets. 18. Why was Fort Margherita built?
  \item For administration
  \item For defense.
  \item For leisure.
\end{choices}
\question{19. What is the distinctive feature of Kuching?}
\begin{choices}(3)
  \item Its cat statues.
  \item Its long history.
  \item Its charming scenery.
\end{choices}
\question{20. What is nasi lemak?}
\begin{choices}(3)
  \item Rice noodles.
  \item Sour soup.
  \item A rice dish.
\end{choices}
\parttitle{第二部分 阅读理解（共两节，满分 50 分）}
\sectitle{第一节（共15小题：每小题2.5分，满分37.5分）}
\par 阅读下列短文，从每题所给的 A、B、C 和 D 四个选项中，选出最佳答案，\par
\passagelabel{A}
\par Reach for the Stars\par
\par Astronomy is a subject that integrates various fields, from physics and mathematics to chemistry and biology. By studying the universe, students gain a comprehensive view of science and see how different fields interconnect. This promotes a deeper understanding of the world around them.\par
\par That's where Space Science Learning Club (SSLC) steps in. SSLC is an after-school astronomy club for school students. At SSLC, we are passionate about inspiring the next generation of space explorers. Our goal is to help children unlock their potential with astronomy and space science, no matter their age or background.\par
\par Led by Learning Coach Vinod Kumar, a NASA certified trainer, our sessions are tailored to each student's individual needs, providing them with the knowledge and skills to awaken their minds. Come and join us on a journey\par
\par to the stars.\par
\par Program Schedule:\par
\par\vspace{0.4em}
\begin{center}\begin{tabular}{{|l|p{\dimexpr\linewidth-3em-4\tabcolsep}|}}
\hline
Week 1 & Naked Eye: Astronomy Night Sky Watch \\
Week 2 & Ancient Astronomy: Tools of Ancient Astronomers \\
Week 3 & Hands-on Astronomy: Create Observation Tools by Yourself \\
Week 4 & Introduction to Stars: Make Your Own Star Chart \\
Week 5 & Digital Astronomy Introduction \\
\hline
\end{tabular}\end{center}
\vspace{0.4em}
\par Each class will be for a period of 3 hours after school hours. 1 class every week.\par
\par Classes are 60\% application and 40\% theory.\par
\par Classes focus on both activity-based and play-driven learning.\par
\par Program Highlights: √ On-site interaction with astronomers √ FREE access to web-based learning √ 10+ hours of classroom training √ On-call coaching assistance for assignments √ Worksheets for every class √ One star observing session through telescope FOR MORE INFORMATION CALL US AT 9360841234\par
\question{21. Which of the following best describes the feature of astronomy?}
\begin{choices}(2)
  \item Culture-specific.
  \item Cross-disciplinary.
  \item Newly-emerging.
  \item Experiment-centered.
\end{choices}
\question{22. In which week can students DIY their star maps?}
\begin{choices}(2)
  \item Week 2.
  \item Week 3.
  \item Week 4.
  \item Week 5.
\end{choices}
\question{23. What does the program provide to students?}
\begin{choices}(2)
  \item Instant coaching support.
  \item Unlimited use of telescope.
  \item Virtual talk with astronomers.
  \item Training on time management.
\end{choices}
\passagelabel{B}
\par During my early years, I focused on eye-catching landscape photography. I spent time in the Himalayas, Patagonia and Tasmania and returned with pictures of breathtaking beauty, which I thought were "awesome". But a one-month camping trip to the Menindee lakes along Darling-Baaka in Australia shifted my perspective: Photography is more than taking photos.\par
\par As I became increasingly familiar with the place, I learned that human activities led to widespread deforestation and drought in the area, resulting in many dead trees. Driving across the region, I recalled that Sydney was blanketed by smoke for months on end as bushfires swept through the east coast. Fortunately, the La Nina weather cycle reached Australia, increasing rainfall and causing the continent's short-lived rivers to spring back to life. Yet, the image of bushfires stayed with me, and I felt driven to do something. That was when the photo Blaze came into being.\par
\par I spoke to the local authority, explaining that I hoped to create a ‘burning’ tree as a beacon(灯塔) to draw people's attention to deforestation. They showed great interest because fire is an integral part of the Australian farming and the natives.\par
\par I connected a few of the gas canisters(罐) generally used for barbecues to flexible gas lines called "slinkies", which could be wired up the back of the tree, away from the tree and out of sight of the camera. These produced an outline of flames(火焰) when I lit them, as if the tree were on fire. It was lit for 15 seconds while I got the shot, and fire-fighting equipment was prepared. Even though the tree was dead, it was still a vital habitat. But the gas flames did no damage and the insects that lived inside were unharmed.\par
\par I love that when people look at the photo, they can concretize the delicacy of our natural world while appreciating Blaze.\par
\question{24. What is the author's new perception of photography?}
\begin{choices}(2)
  \item It involves professional skills.
  \item It should carry deep meanings.
  \item It requires high-end equipment.
  \item It should produce beautiful pictures.
\end{choices}
\question{25. What inspired the author to create Blaze?}
\begin{choices}(2)
  \item Local environmental changes.
  \item Traditional Australian farming.
  \item Scenery of the Menindee lakes.
  \item Support from the local authority.
\end{choices}
\par 26. Which of the following images best illustrates the photo Blaze?\par
\par [图片] A\par
\par [图片] B\par
\par [图片] C\par
\par [图片] D\par
\question{27. What is the photo Blaze intended to be?}
\begin{choices}(2)
  \item A means of enhancing artistic taste.
  \item A strategy to beautify the landscape.
  \item A call to raise public awareness.
  \item A window to showcase local culture.
\end{choices}
\passagelabel{C}
\par Numerous industries produce larger quantities of wastewaters containing high concentrations of organic substances and inorganic salts. Discharging these wastewaters to the environment is harmful to ecosystem and human health. However, the high salt content in these wastewaters pose serious challenges for existing treatment processes. Researchers at Rice University, in partnership with Guangdong University of Technology, have developed a groundbreaking method for it—adopting dialysis (透析) technology commonly used in medicine.\par
\par The researchers selected commercial ultrafiltration membranes(超虑膜) which were designed to let salt pass through while trapping organic substances. They set up a system where salty waste water and fresh water flowed in opposite directions on either side of a special membrane, without needing pumps. Due to the difference of salt concentrations between waste water and fresh water, salt naturally moves into the fresh water while organic substances are kept back. To test this, they compared the amount of organic material and salt in the\par
\par [图片]\par
\par waste water before and after treatment. The results showed the system could effectively remove salt without using extra amounts of fresh water.\par
\par One significant advantage of dialysis is its resistance to fouling (积垢). Unlike pressure-driven systems, dialysis experienced notably less buildup of organic materials on the membrane because it doesn't rely on extra pressure. "This could translate to lower energy consumption, less maintenance, and fewer membrane replacements," said Elimelech, a researcher on the study.\par
\par Additionally, with suitable membranes, dialysis can be tailored to separate various organics, enabling the selective recovery of valuable resources from the wastewater. In this way, dialysis creates opportunities for transforming wastewater treatment plants into resource recovery factories, aligning the principles of a recycling economy.\par
\par Overall, while dialysis alone doesn't fully purify waste water, it can transform the high-salinity organic wastewater stream into a low-salinity organic wastewater stream that can be readily treated by conventional methods, which revolutionizes the treatment of some of the most challenging industrial wastewaters.\par
\question{28. What is the major challenge faced by current wastewater treatment?}
\begin{choices}(2)
  \item The lack of sufficient treatment devices.
  \item The high salt percentage in wastewaters.
  \item The substantial volumes of wastewaters.
  \item The considerable cost of treatment processes.
\end{choices}
\question{29. What can be learned about the process of dialysing wastewaters?}
\begin{choices}(2)
  \item The membrane transformed salt into organics.
  \item The extra pressure contributed to salt transport.
  \item The concentration difference played the key role.
  \item Liquids on both sides flowed in the same direction.
\end{choices}
\question{30. What does the underlined word "aligning" in paragraph 4 probably mean?}
\begin{choices}(2)
  \item Matching.
  \item Establishing.
  \item Adjusting.
  \item Challenging.
\end{choices}
\question{31. Which of the following can be a suitable title?}
\begin{choices}(2)
  \item Dialysis: A Pioneering Blessing for Medicine
  \item Membranes: A Protective Screen for Fresh Water
  \item Dialysis: A New Pathway for Wastewater Treatment
  \item High-Salinity Wastewater: A Threat to the Environment.
\end{choices}
\passagelabel{D}
\par At first glance, the scientist Charles Darwin and the poet Emily Dickinson would seem to have little in common. In Natural Magic, literary scholar Renée Bergland seeks to convince the reader otherwise, interweaving the lives, work, and history of these two towering figures of the 19th century. In doing so, Bergland shows definitively that they shared deep sensibilities about the natural world. They were both keen observers in their own home gardens. For both, the "magic" of nature could be found in the physical and material rather than the supernatural. They found nature to be a place enchanted by its own mysteries, joys, and sorrows.\par
\par Taking the form of a joint biography(传记), Natural Magic alternates between Darwin and Dickinson over the course of 16 chapters, progressing simultaneously(同时地) in time. While Bergland offers comprehensive descriptions, building on the extensive work of other biographers and scholars, the book's own magic shines in the conversation created between its subjects' bodies of work. When viewed as part of a historical dialogue, Darwin and Dickinson amplify(阐发) each other's work, revealing the science in her poetry and the poetry in his science.\par
\par Natural Magic further situates both writers in the history of science as it developed from "natural philosophy" into specialized disciplines throughout the 19th century. Bergland lays extensive groundwork for the social and cultural forces that shaped each thinker in a time when the sciences and arts were more closely connected.\par
\par Bergland herself displays moments of beauty in her writing, particularly in her vivid descriptions of the wonder of scientific discoveries in the 19th century. "Changing angles of sunlight, the unchanging stars, and . the bright colors of the great auroras lightened up the vastness of planetary space," she writes, for example.\par
\par Ultimately, the book offers readers a historical and biographical look into the magic offered by the natural world as reflected in the works of two major 19th-century thinkers. Bergland concludes that in a time of climate crisis, we must recover this kind of dialogue across the sciences and the arts, emphasizing our shared sense of wonder in nature.\par
\question{32. What does Bergland aim to demonstrate about Darwin and Dickinson?}
\begin{choices}(2)
  \item Their inborn talents.
  \item Their great achievements.
  \item Their literary background.
  \item Their unseen commonalities.
\end{choices}
\question{33. What can be learned concerning Natural Magic?}
\begin{choices}(2)
  \item It creates dialogues between works.
  \item It is a co-authored masterpiece.
  \item It contradicts other scholars' views.
  \item It focuses on natural philosophy.
\end{choices}
\question{34. Which best describes Bergland's language in describing the wonder of scientific discoveries?}
\begin{choices}(2)
  \item Straightforward.
  \item Academic.
  \item Humorous.
  \item Poetic.
\end{choices}
\question{35. Which of the following might Bergland agree with?}
\begin{choices}(2)
  \item History should never be forgotten.
  \item Nature should be treated with awe.
  \item Wonder is the beginning of wisdom.
  \item The dawn of science came too late.
\end{choices}
\sectitle{第二节（共5小题；每小题2.5分，满分12.5分）}
\par 阅读下面短文，从短文后的选项中选出能填入空白处的最佳选项。选项中有两项为多余选项，\par
\par Creating your own plant specimens(标本) is a fun, simple project that anyone can do at home. It involves collecting plant specimens, pressing them and labeling them. \blanknum{36} Just keep an eye out wherever you go and look for your eye-catching plants. Labeling them is a topic I'll save for another day. So today I want to show you how to press the specimens.\par
\par Construct a homemade plant press. \blanknum{37} Then cover it with a piece of corrugated (起皱的) card, followed by 2 sheets of absorbent paper. This will serve as the bottom half of your press. Have a second board and more cards and absorbent paper on standby for when you're ready to put the top half.\par
\par Arrange your specimens on the bottom half of your press. Lay out each cutting which you've cleaned beforehand in such a way that it looks as attractive as possible. \blanknum{38} If possible, stretch smaller specimens out to their full length so that they dry perfectly flat. Leave 1—2 inches of space between each specimen to keep them from sticking together.\par
\par \_39 Once you're satisfied with the layout of your plants, sandwich them with another layer of absorbent paper, corrugated card, and hardboard. Pile bricks, large books, or similar items evenly across the upper board to weigh it down and press the cuttings flat.\par
\par Continue drying and pressing your specimens. \blanknum{40} Depending on the size and type of the plants you're preserving, it may take them anywhere from a couple of days to 3 full weeks to dry completely.\par
\begin{choices}(1)
  \item Collecting plant specimens is very simple.
  \item Describe any outstanding features you notice.
  \item Start by selecting the materials you may need.
  \item Place one hardboard flat against a level surface.
  \item The duration of this process varies considerably.
  \item Close the press and set heavy objects on top to create pressure.
  \item Meanwhile, ensure each of its main features can be seen clearly.
\end{choices}
\parttitle{第三部分 语言知识运用（共两节，满分 30 分）}
\sectitle{第一节（共15小题；每小题1分，满分15分）}
\par 阅读下面短文，从短文后各题所给的 A、B、C 和 D 四个选项中，选出可以填入空白处的最佳选项，I never knew a chance encounter was to change my career forever.\par
\par One snowy December day, I wandered out of my New York apartment and \blanknum{41} for a grocery store on Third Avenue. Behind me someone shouted: "Where can I find peaches?" I turned around to \blanknum{42} an annoyed woman standing outside the greengrocer's. The absurdity(荒谬) of the moment \blanknum{43} me—why must someone\par
\par buy peaches in the middle of winter?\par
\par I was already \blanknum{44} of the issues facing the food system: industrialized farming destroying our soils, unclear supply chains leaving citizens \blanknum{45} in making the right buying decisions, and the \blanknum{46} of ultra-processed foods with zero nutritional value in supermarkets, schools and hospitals, to name a few. But this moment underscored our serious \blanknum{47} with nature and its seasons. We had \blanknum{48} the idea that food could and should be eaten any time of the year. Too terrible! An inner voice whispered: "Do something to \blanknum{49} seasonal flavour into people's lives."\par
\par After graduation, \blanknum{50} by the skyscrapers(摩天大楼) and energy of the Big Apple, I landed a decent job with a \blanknum{51} salary, managing financial investment for clients. Yet my work was stuck in an undercurrent of greed that I struggled to \blanknum{52}. Now it was time.\par
\par I was \blanknum{53} my 40s. If I did not act now, it would be too late. An urgency \blanknum{54} and food became my "North Star". That was the launcher that \blanknum{55} me from a New York skyscraper to a London food market.\par
\cz{41}{spoke}{cheered}{worked}{headed}
\cz{42}{spot}{hit}{approach}{comfort}
\cz{43}{amused}{tricked}{struck}{hurt}
\cz{44}{afraid}{aware}{ashamed}{careful}
\cz{45}{impatient}{reasonable}{picky}{powerless}
\cz{46}{assessment}{prejudice}{dominance}{classification}
\cz{47}{dissatisfaction}{disconnect}{familiarity}{mixture}
\cz{48}{normalized}{finalized}{commercialized}{prioritized}
\cz{49}{program}{fold}{channel}{transform}
\cz{50}{greeted}{attracted}{surrounded}{refreshed}
\cz{51}{flexible}{basic}{handsome}{fair}
\cz{52}{escape}{maintain}{interpret}{reveal}
\cz{53}{averaging}{nearing}{braving}{battling}
\cz{54}{kicked in}{built up}{passed by}{caught on}
\cz{55}{witnessed}{accompanied}{saved}{shot}
\sectitle{第二节（共10小题；每小题1.5分，满分15分）}
\par 阅读下面短文，在空白处填入 1 个适当的单词或括号内单词的正确形式，\par
\par The recent announcement by China's National Health Commission (NHC) \blanknum{56} it would promote a three-year "weight management" campaign has fueled nationwide enthusiasm. This campaign moves beyond traditional policy approaches \blanknum{57} combining nationwide cooperation with community-level engagement.\par
\par Over past years, China has witnessed an upward trend in the disease rate of its overweight and obese population. Obesity (肥胖) has become a major public health issue in China, ranking as the \blanknum{} (five) leading risk factor for death and disability in the country.\par
\par What \blanknum{59} (set) this campaign apart is its comprehensive strategy, excelling through cross-sector cooperation. Sixteen government agencies \blanknum{60} (joint) work to address obesity's complex challenges, creating targeted plans for different age groups and life phases. \blanknum{61} (advance) the mission, the NHC crafted the "Nutrition \& Exercise Guide", which has gone viral for its localized eating advice that harmonizes cultural food traditions with science-backed \blanknum{62} (well) principles.\par
\par Acknowledging obesity's physical and emotional impacts, the campaign provides personalized health blueprints \blanknum{63} (feature) meal schedules and workout routines. Furthermore, specialized teams composed of dietitians, traditional medicine experts, and fitness trainers create personalized support networks, monitoring daily progress via health apps, transforming health awareness into \blanknum{} \blanknum{64}\blanknum{} engaging daily habit.\par
\par\vspace{0.4em}
\begin{center}\begin{tabular}{{|l|}}
\hline
Winners Never Quit \\
\_ \\
\hline
\end{tabular}\end{center}
\vspace{0.4em}
\par It's hoped that strong centralized leadership \blanknum{65} (pair) with local engagement can effectively address obesity.\par
\parttitle{第四部分 写作（共两节，满分 40 分）}
\sectitle{第一节（满分15分）}
\par 你校英文报 Youth Voice 栏目正在举办征文比赛。请你以 "Winners never quit" 为题写一篇英文短文参赛，内容包括：(1) 对该格言的理解；(2) 个人经历及感悟。\par
\par 注意：(1) 写作词数应为 80 左右；(2) 请按如下格式在答题纸的相应位置作答，\par
\sectitle{第二节（满分25分）}
\par 阅读下面材料，根据其内容和所给段落开头语续写两段，使之构成一篇完整的短文，\par
\par My two best friends were each doing something exciting — Carlos was at the beach rescuing sea turtles, and Natalia was with her mother, a famous designer working for a celebrity in a big city. But for me, things were quite the opposite.\par
\par I complained to Mom over the phone about the extreme boredom during the first week of my summer stay on my grandparents' farm. She said, "Don't envy others! Just dive into the experience and taste true joy. Not everyone gets the opportunity to live on a farm."\par
\par Her words gave me an idea, so I decided I would compose my social media status updates on the farm. Then, my post hit the feed the next morning:\par
\par Alyssa's Status:5:54 A.M. The rooster(公鸡) is standing tall on his post, wings shining in the dawn light. He beats his wings, lifts his héd high, and declares the start of a brand new day.\par
\par Carlos posted at 7:39 A. M., saying he was searching for sea turtle nests on the beach and hoping to find one. Natalia posted at 11:19 A. M., asking which color to paint a rock star's house. They got 26 and 28 likes respectively. To my great joy, my post got almost the same amount of attention as theirs. Satisfying!\par
\par Just as I was about to comment on their posts, Grandma interrupted me, "Alyssa, put your shoes on. We're off to Haverford's farm to pick blackberries." Later that afternoon, I posted about the unexpected fun and described in vivid detail about the fantastic flavors of blackberry pies.\par
\par Responses? Carlos said he wished he could taste one while Natalia was envious of my farm adventure. Interesting!\par
\par After breakfast the following morning, Grandma said, "I was just thinking that today might be a good day for you to explore the old barn (谷仓). I’ve spotted foxes nearby—maybe you’ll find some adorable baby foxes there."\par
\par "That would have potential for a worthy photo to post online," I thought to myself. Perfect! With curiosity, I approached the old barn, and pushed the door open.\par
\par 注意：(1) 续写词数应为 150 左右；(2) 请按如下格式在答题卡的相应位置作答。\par
\writingpara{Paragraph 1: Suddenly, a weak sound was heard from a corner of the barn.}
\writinglines{8}
\writingpara{Paragraph 2: As I was waiting for the comments, something within me started to change.}
\writinglines{8}
\end{document}
